\documentclass[11pt,oneside]{book}
\usepackage[paperwidth=6.5in,paperheight=9.25in,top=0.82in,bottom=0.8in,inner=0.9in,outer=0.72in,headheight=14pt]{geometry}
\usepackage[T1]{fontenc}
\usepackage{lmodern}
\usepackage{microtype}
\usepackage{xcolor}
\usepackage{booktabs}
\usepackage{hyperref}
\definecolor{openblue}{HTML}{1E5A85}
\hypersetup{colorlinks=true,linkcolor=openblue,citecolor=openblue,urlcolor=openblue}
\title{Open Tools, Open Books\\[0.4em]\large A bookdown-inspired textbook starter}
\author{Your Name}
\date{\today}
\begin{document}
\frontmatter
\maketitle
\chapter*{How to use this book}
This starter is a small, self-contained example of the chapter structure used by an open textbook workflow. Replace the chapters with your own source, then choose whether your publishing pipeline should produce a web book, PDF, or EPUB.
\tableofcontents
\mainmatter
\chapter{Open authoring}
An open textbook is easier to maintain when its source, data, figures, and build instructions travel together. Keep a short overview at the front of the project and make each chapter useful on its own.
\section{A reproducible chapter}
Write the learning goal first. Then add an example, a short activity, and a recap. The same source can be rendered into different formats when structure is kept separate from presentation.
\begin{table}[ht]
\centering
\caption{A small chapter plan}
\begin{tabular}{lll}
\toprule
Part & Purpose & Evidence \\
\midrule
Overview & Orient the reader & One-sentence goal \\
Example & Show the idea & Reproducible result \\
Activity & Invite practice & Reader response \\
Recap & Reinforce the idea & Three key points \\
\bottomrule
\end{tabular}
\end{table}
\chapter{Tools and evidence}
Document the tools used to create figures, analyses, and references. A reader should be able to distinguish a claim, its evidence, and the version of the tool that produced it.
\section{Publishing paths}
The upstream OER bookdown project demonstrates web, PDF, and EPUB output from one open textbook source. This LaTeX-only starter keeps the PDF path immediately runnable while leaving room to connect a bookdown or Quarto build later.
\chapter{Revision and sharing}
Invite readers to report errors and unclear explanations. A changelog and a stable citation make improvements visible without hiding the history of the book.
\backmatter
\chapter*{References and further reading}
Add references here or replace this chapter with a generated bibliography from your preferred toolchain.
\end{document}
